\documentclass[10pt,a4paper]{amsart}
\usepackage[margin=19mm]{geometry}
\usepackage{amsmath,amssymb,mathtools}
\usepackage[scr=rsfs,cal=euler]{mathalfa}
\usepackage{xfrac}
\usepackage[unicode,hidelinks,bookmarks=false]{hyperref}
\newcommand{\ie}{i.e.\ }
\newcommand{\cf}{cf.\ }
\newcommand{\resp}{respectively\ }
\newcommand{\Subsection}{\subsection}
\providecommand{\nobreakdash}{\nobreak\hbox{-}}
\setlength{\emergencystretch}{2em}
\multlinegap0pt
\usepackage{bm}
\usepackage{bbm}
\usepackage{dsfont}
\usepackage{xspace}
\usepackage{cancel}
\usepackage{mathtools}
\usepackage{amsthm}
\makeatletter
\@ifundefined{thm}{\newtheorem{thm}{Thm}}{}
\@ifundefined{lem}{\newtheorem{lem}[thm]{Lemma}}{}
\makeatother
\makeatletter
\expandafter\ifx\csname AMS\endcsname\relax
\newenvironment{AMS}{}{}
\fi
\DeclareMathOperator{\N}{\mathbb{N}}
\DeclareMathOperator{\Q}{\mathbb{Q}}
\DeclareMathOperator{\Z}{\mathbb{Z}}
\expandafter\ifx\csname keywords\endcsname\relax
\newenvironment{keywords}{}{}
\fi
\DeclareMathOperator{\supp}{{\rm Supp}}
\makeatother
\title[Cyclotomic factors of rational necklace functions]{Cyclotomic factors of rational necklace functions}
\author{Nguyen Cao Minh, Nguyen Vu Hoang Minh, Dung Nguyen, Tung T. Nguyen, Nguyen Duy Tan, Duong Tran}
\date{Source: arXiv:2606.02324v1 (CC BY 4.0)}
\begin{document}
\maketitle

\noindent\emph{Source and licence.} Cyclotomic factors of rational necklace functions. Version arXiv:2606.02324v1 (CC BY 4.0), distributed under the Creative Commons Attribution 4.0 International licence, as recorded in the arXiv licence metadata for this exact version and shown on the abstract page https://arxiv.org/abs/2606.02324v1. Full rights notice: licenses/arxiv-2606.02324-CC-BY-4.0.txt.\par\smallskip

\noindent\textit{Editorial scope.} Counted unit: the Lem whose source label is \texttt{None} in the article (source0.tex lines 1574--1576, source label \texttt{None}), delivered together with the article's own complete original proof of it (source0.tex lines 1577--1670, one \texttt{proof} environment of 94 lines). No mathematics is written, repaired or re-derived: statement and proof are the article's own text line for line. Required context retained uncounted: none. Every definition, display and label of those lines is retained unchanged. Omitted from this package and not redistributed: 1--1573, 1671--1916. Nothing inside a retained excerpt is omitted or altered: every retained body is delivered exactly as the article's lines are written, so no omission receipt of any kind accompanies this package. Citations: the retained excerpts carry no citation command at all, so no bibliography entry of the article is transcribed and no reference is redistributed. Reference adaptation: none was needed and none was made; every reference inside the delivered unit resolves inside it, so no unresolved reference or citation is produced. Printed slips kept literal: no printed word, symbol, hypothesis or number of the counted statement or of its proof was altered. Layout and class: the article's own class is \texttt{amsart} with packages including amsmath, amssymb, amsfonts, float, xy, caption, enumerate, mathpazo, a4wide, ulem, bm, graphicx, hyperref, xcolor; the wrapper is the lane's pinned \texttt{amsart} editorial wrapper at 10pt on A4, which restates the theorem-like environments the retained text needs and adds the article's own macro definitions that the retained text needs (\texttt{\textbackslash N}, \texttt{\textbackslash Q}, \texttt{\textbackslash Z}, \texttt{\textbackslash supp}), with extra wrapper packages (bm, bbm, dsfont, xspace, cancel, mathtools). The delivered numbering is the unit's own. Assets: no retained span contains a figure environment, a graphics-inclusion command, a PDF-page inclusion, a table or a TikZ picture; no image file is copied into this package, the package declares no asset dependency and no image file is redistributed with it. Licence: version arXiv:2606.02324v1 of this article is distributed under the Creative Commons Attribution 4.0 International licence, as recorded in the arXiv licence metadata for this exact version and shown on the abstract page https://arxiv.org/abs/2606.02324v1. A full rights notice is delivered at licenses/arxiv-2606.02324-CC-BY-4.0.txt. Characters: every retained excerpt is pure ASCII, so the delivered engine, XeLaTeX with its default Latin Modern text font, reports no missing character.
\begin{lem} Let $\alpha\in \Z[\N^\circ]$ and $f\in \Q(x)\setminus \Q$. If $\alpha f=0$ then  $\alpha=0$.
    
\end{lem}

\begin{proof}
 Write
\( 
\alpha=\sum_{m\in S} a_m[m],
\)
where \(S=\supp(\alpha)\subset \mathbb N^\circ\) is finite and \(a_m\neq 0\) for all \(m\in S\).
We will show that \(\alpha f=0\) implies \(S=\varnothing\), hence \(\alpha=0\).

Write \(f=P/Q\), where \(P,Q\in \mathbb Q[x]\) are coprime and \(Q\neq 0\). Since \(f\notin \mathbb Q\), at least one of \(P,Q\) has positive degree. 

Suppose
\[
0=\alpha f(x)=\sum_{m\in S} a_m f(x^m)=\sum_{m\in S} a_m \frac{P(x^m)}{Q(x^m)}.
\]
Multiplying by
\( 
\prod\limits_{m\in S} Q(x^m),
\)
we get
\[
\sum_{m\in S}
a_m P(x^m)\prod_{\substack{n\in S\\ n\neq m}} Q(x^n)=0.
\]

 The summand corresponding to  \(m\) has a degree
\[
d_m:=m\deg P+\sum_{\substack{n\in S\\ n\neq m}} n\deg Q =\sum_{n\in S} n\deg Q + m(\deg P-\deg Q).
\]
Hence for $m,m'\in S$, we have
\[
d_m-d_{m'}=(m-m')(\deg P-\deg Q). 
\]
Suppose $\deg P\not=\deg Q$ then all  degrees $d_m$ are distinct. This forces $S$ is empty and $\alpha=0$.

 Assume now that \(\deg P=\deg Q\). Let
\(
c=\dfrac{\operatorname{lc}(P)}{\operatorname{lc}(Q)}
\)
and set
\(
R=P-cQ.
\) 
Then
\(
f=c+\frac{R}{Q}=c+g,
\)
and
\(
\deg R<\deg Q.
\).

Since \(\alpha f=0\), we have
\( 
\alpha\!\left(\frac{R}{Q}\right)
=
-c\,\alpha(1).
\)
The right-hand side is a constant, denoted $C\in \Q$. We obtain
\[
\sum_{m\in S} a_m \frac{R(x^m)}{Q(x^m)}=C.
\]


Multiplying by \(\prod_{m\in S}Q(x^m)\) yields
\[
\sum_{m\in S}
a_m R(x^m)
\prod_{\substack{n\in S\\ n\neq m}}
Q(x^n)
=
C\prod_{m\in S}Q(x^m).
\]

If $C\not=0$ then the degree of the right-hand side is
\(
\sum\limits_{m\in S} m\deg Q.
\)
On the other hand, each summand on the left has degree
\[
m\deg R+
\sum_{\substack{n\in S\\ n\neq m}}
n\deg Q
=
\sum_{n\in S} n\deg Q
-
m(\deg Q-\deg R).
\]
Since \(\deg R<\deg Q\), every summand on the left has degree strictly less than
\(\sum_{n\in S} n\deg Q\). Therefore \(C=0\), and hence
\(
\alpha g=0
\), where $g=\dfrac{R}{Q}$ with \(\deg R<\deg Q\). Note that  since $f$ is not a constant $g$ is also not a constant.
Thus we reduce to the previous case and hence  $\alpha=0$.
\end{proof}

\end{document}
